\begin{exercise}{27}
Let $f\colon X\to Y$ be a map of simplicial sets. Prove or give a counter example to
the following statements:
\begin{enumerate}[label=(\arabic*),leftmargin=*,itemsep=0pt,parsep=0pt,
                  topsep=.5\baselineskip]
 \item If $f$ is a monomorphism, then $hX\to hY$ is fully faithful.
 \item If $f$ is a degree-wise surjection, then $hX\to hY$ is
       \emph{surjective and full}, i.e., it induces a surjection on objects and on hom-sets.
 \item If $f$ induces a surjection on $0$- and $1$-simplices, then
       $hX\to hY$ is surjective and full.
\end{enumerate}

\begin{proof}
All three statements are false.

(1) Take the inclusion $\Lambda_0^2\hookrightarrow\Delta^2$.
By Exercise~\ref{ex:26}, the source homotopy category is the span
$1\leftarrow0\to2$, whereas $h(\Delta^2)=[2]$.
Thus the hom-map from $1$ to $2$ is $\varnothing\to\{1\to2\}$,
which is not surjective.

(2), (3) Let
\[
 X=\coprod_{0\le i<j\le2}\Delta^1
 \quad\text{and}\quad
 p\colon X\longrightarrow\partial\Delta^2
\]
map each summand onto the edge $(i,j)$.
Every simplex of $\partial\Delta^2$ has image in an edge, so
$p$ is degree-wise surjective and satisfies both hypotheses.
By Definition~1.2.5, $hX$ is the disjoint union of three copies
of $[1]$: no edge, and hence no path, joins different summands.
The initial vertex of the $(0,1)$ summand and the terminal
vertex of the $(1,2)$ summand have no morphism between them, whereas
their images do in $h(\partial\Delta^2)$. Thus $h(p)$ is not full.
In fact, its image contains only identities and the three
generating edges; the distinct composite $0\to1\to2$ in the
free category $h(\partial\Delta^2)$ of Exercise~\ref{ex:26}
is not in the image either.
\end{proof}
\end{exercise}
